% GENERATED FILE. Edit canonical lessons, then run npm run generate:books.
% canonical-source-hash: fnv1a64:8cdcb0d33f2fcfd5
% canonical-lessons: GE-C509-wettkampf, GE-C509-meisterschaft, GE-C509-pokal, GE-C509-medaille, GE-C509-siegerehrung

\chapter{Pokal und Medaille}
\label{ch:ge-pokal-und-medaille}
\hlchaptermodality{509}

\begin{chapteropening}
\textbf{By the end of this chapter:} \emph{I can talk about competitions and prizes.}

Complete the last lesson of chapter 509: I can talk about competitions and prizes.
\end{chapteropening}

\section[der Wettkampf]{der Wettkampf — competition, contest}

\label{lesson:GE-C509-wettkampf}



\begin{quote}
Before the new one: say the German for opponent, then the German for tournament.
\end{quote}

\subsection*{You'll want to know: der Wettkampf}
\textbf{der Wettkampf} — \textquotedblleft{}competition, contest\textquotedblright{}.

\begin{quote}
Beim Wettkampf war sie sehr nervös.
\end{quote}

\begin{quote}
\emph{She was very nervous at the competition.}
\end{quote}

Plural: \emph{die Wettkämpfe}. For sport; a general contest is \emph{der Wettbewerb}.

\begin{cousinweb}
\emph{Wettkampf} is built on \emph{Wette} \textquotedblleft{}bet\textquotedblright{}, inherited from Old High German \emph{wetti}, and \emph{Kampf} \textquotedblleft{}fight\textquotedblright{}, which came through Old High German \emph{kampf} from Latin \emph{campus} \textquotedblleft{}field\textquotedblright{} (EtymDB).
\end{cousinweb}

\begin{grammarlens}[title={What you've built}]
One more word for telling, explaining and describing at B1.
\end{grammarlens}

\subsection*{Your turn}
\begin{itemize}
\raggedright
  \item \emph{Say it:} \emph{der Wettkampf}
  \item \emph{Say it:} \emph{der Wettkampf}, once more
  \item \emph{Say it:} a sentence of your own with \emph{der Wettkampf}
  \item \emph{Recall:} say the German for opponent, then the German for tournament, then say \emph{der Wettkampf} again
\end{itemize}

\begin{tcolorbox}[breakable,colback=teal!4,colframe=teal!35!black,title={Before you move on}]
What does der Wettkampf mean? (\textquotedblleft{}Competition, contest\textquotedblright{}.) Say it once more.
\end{tcolorbox}
\section[die Meisterschaft]{die Meisterschaft — championship}

\label{lesson:GE-C509-meisterschaft}



\begin{quote}
Before the new one: say the German for tournament, then the German for competition.
\end{quote}

\subsection*{You'll want to know: die Meisterschaft}
\textbf{die Meisterschaft} — \textquotedblleft{}championship\textquotedblright{}.

\begin{quote}
Unser Verein hat die Meisterschaft gewonnen.
\end{quote}

\begin{quote}
\emph{Our club won the championship.}
\end{quote}

Plural: \emph{die Meisterschaften}. \emph{Die Weltmeisterschaft} — the World Cup, a world championship.

\begin{cousinweb}
\emph{Meisterschaft} is built on \emph{Meister} \textquotedblleft{}master, champion\textquotedblright{}, which came through Old High German \emph{meistar} from Latin \emph{magister} (EtymDB), with \emph{-schaft}.
\end{cousinweb}

\begin{grammarlens}[title={What you've built}]
One more word for telling, explaining and describing at B1.
\end{grammarlens}

\subsection*{Your turn}
\begin{itemize}
\raggedright
  \item \emph{Say it:} \emph{die Meisterschaft}
  \item \emph{Say it:} \emph{die Meisterschaft}, once more
  \item \emph{Say it:} a sentence of your own with \emph{die Meisterschaft}
  \item \emph{Recall:} say the German for tournament, then the German for competition, then say \emph{die Meisterschaft} again
\end{itemize}

\begin{tcolorbox}[breakable,colback=teal!4,colframe=teal!35!black,title={Before you move on}]
What does die Meisterschaft mean? (\textquotedblleft{}Championship\textquotedblright{}.) Say it once more.
\end{tcolorbox}
\section[der Pokal]{der Pokal — trophy, cup}

\label{lesson:GE-C509-pokal}



\begin{quote}
Before the new one: say the German for competition, then the German for championship.
\end{quote}

\subsection*{You'll want to know: der Pokal}
\textbf{der Pokal} — \textquotedblleft{}trophy, cup\textquotedblright{}.

\begin{quote}
Der Kapitän hob den Pokal in die Luft.
\end{quote}

\begin{quote}
\emph{The captain lifted the trophy into the air.}
\end{quote}

Plural: \emph{die Pokale}.

\begin{cousinweb}
\emph{Pokal} was borrowed from Italian \emph{boccale} \textquotedblleft{}jug\textquotedblright{} (EtymDB), from Late Latin \emph{baucalis}, from Greek: a drinking cup before it was a prize.
\end{cousinweb}

\begin{grammarlens}[title={What you've built}]
One more word for telling, explaining and describing at B1.
\end{grammarlens}

\subsection*{Your turn}
\begin{itemize}
\raggedright
  \item \emph{Say it:} \emph{der Pokal}
  \item \emph{Say it:} \emph{der Pokal}, once more
  \item \emph{Say it:} a sentence of your own with \emph{der Pokal}
  \item \emph{Recall:} say the German for competition, then the German for championship, then say \emph{der Pokal} again
\end{itemize}

\begin{tcolorbox}[breakable,colback=teal!4,colframe=teal!35!black,title={Before you move on}]
What does der Pokal mean? (\textquotedblleft{}Trophy, cup\textquotedblright{}.) Say it once more.
\end{tcolorbox}
\section[die Medaille]{die Medaille — medal}

\label{lesson:GE-C509-medaille}



\begin{quote}
Before the new one: say the German for championship, then the German for trophy.
\end{quote}

\subsection*{You'll want to know: die Medaille}
\textbf{die Medaille} — \textquotedblleft{}medal\textquotedblright{}.

\begin{quote}
Sie gewann eine Medaille im Weitsprung.
\end{quote}

\begin{quote}
\emph{She won a medal in the long jump.}
\end{quote}

Plural: \emph{die Medaillen}. Said roughly \emph{me-DAL-ye}.

\begin{cousinweb}
\emph{Medaille} came from French \emph{médaille}, from Italian \emph{medaglia} (EtymDB), which goes back to Latin \emph{metallum} \textquotedblleft{}metal\textquotedblright{}.
\end{cousinweb}

\begin{grammarlens}[title={What you've built}]
One more word for telling, explaining and describing at B1.
\end{grammarlens}

\subsection*{Your turn}
\begin{itemize}
\raggedright
  \item \emph{Say it:} \emph{die Medaille}
  \item \emph{Say it:} \emph{die Medaille}, once more
  \item \emph{Say it:} a sentence of your own with \emph{die Medaille}
  \item \emph{Recall:} say the German for championship, then the German for trophy, then say \emph{die Medaille} again
\end{itemize}

\begin{tcolorbox}[breakable,colback=teal!4,colframe=teal!35!black,title={Before you move on}]
What does die Medaille mean? (\textquotedblleft{}Medal\textquotedblright{}.) Say it once more.
\end{tcolorbox}
\section[die Siegerehrung]{die Siegerehrung — award ceremony}

\label{lesson:GE-C509-siegerehrung}



\begin{quote}
Before the new one: say the German for trophy, then the German for medal.
\end{quote}

\subsection*{You'll want to know: die Siegerehrung}
\textbf{die Siegerehrung} — \textquotedblleft{}award ceremony\textquotedblright{}.

\begin{quote}
Nach dem Rennen gab es eine Siegerehrung auf dem Marktplatz.
\end{quote}

\begin{quote}
\emph{After the race there was an award ceremony in the market square.}
\end{quote}

Plural: \emph{die Siegerehrungen}. \emph{Die Ehrung} — honouring.

\begin{cousinweb}
\emph{Siegerehrung} is built on \emph{Sieger} \textquotedblleft{}winner\textquotedblright{}, from \emph{Sieg}, Proto-Germanic \emph{segaz} (EtymDB), and \emph{Ehrung} \textquotedblleft{}honouring\textquotedblright{}, from \emph{ehren} \textquotedblleft{}to honour\textquotedblright{}.
\end{cousinweb}

\begin{grammarlens}[title={What you've built}]
One more word for telling, explaining and describing at B1.
\end{grammarlens}

\subsection*{Your turn}
\begin{itemize}
\raggedright
  \item \emph{Say it:} \emph{die Siegerehrung}
  \item \emph{Say it:} \emph{die Siegerehrung}, once more
  \item \emph{Say it:} a sentence of your own with \emph{die Siegerehrung}
  \item \emph{Recall:} say the German for trophy, then the German for medal, then say \emph{die Siegerehrung} again
\end{itemize}

\begin{tcolorbox}[breakable,colback=teal!4,colframe=teal!35!black,title={Before you move on}]
What does die Siegerehrung mean? (\textquotedblleft{}Award ceremony\textquotedblright{}.) Say it once more.
\end{tcolorbox}
